
\documentclass{article}
\usepackage{colortbl}
\usepackage{makecell}
\usepackage{multirow}
\usepackage{supertabular}

\begin{document}

\newcounter{utterance}

\centering \large Interaction Transcript for game `cladder', experiment `full\_v1.5\_default', episode 3598 with Qwen3.5{-}27B.
\vspace{24pt}

{ \footnotesize  \setcounter{utterance}{1}
\setlength{\tabcolsep}{0pt}
\begin{supertabular}{c@{$\;$}|p{.15\linewidth}@{}p{.15\linewidth}p{.15\linewidth}p{.15\linewidth}p{.15\linewidth}p{.15\linewidth}}
    \# & \multicolumn{2}{c}{Player} && \multicolumn{2}{c}{Game Master} \\
    \hline

    \theutterance \stepcounter{utterance}  
    & & & \multicolumn{4}{p{0.6\linewidth}}{
        \cellcolor[rgb]{0.9,0.9,0.9}{
            \makecell[{{p{\linewidth}}}]{
                \texttt{\tiny{[P1$\langle$GM]}}
                \texttt{You are an expert in causal inference. The following question is not a typical commonsense query, but rather a meticulously designed question created by a professor specializing in causal inference, intended to assess the students' mastery of the course content.} \\
\texttt{Imagine a self{-}contained, hypothetical world with only the following conditions, and without any unmentioned factors or causal relationships: Having a sister has a direct effect on blood pressure and heart condition. Blood pressure has a direct effect on heart condition. For people who do not have a sister and with low blood pressure, the probability of healthy heart is 63\%. For people who do not have a sister and with high blood pressure, the probability of healthy heart is 31\%. For people who have a sister and with low blood pressure, the probability of healthy heart is 72\%. For people who have a sister and with high blood pressure, the probability of healthy heart is 41\%. For people who do not have a sister, the probability of high blood pressure is 66\%. For people who have a sister, the probability of high blood pressure is 38\%. Does having a sister positively affect heart condition through blood pressure?} \\
\texttt{Start your answer with 'Yes' or 'No', followed by additional reasoning or evidence to support your explanation.} \\
            }
        }
    }
     \\ \\

    \theutterance \stepcounter{utterance}  
    & \multicolumn{4}{p{0.6\linewidth}}{
        \cellcolor[rgb]{0.561,0.737,0.561}{
            \makecell[{{p{\linewidth}}}]{
                \texttt{\tiny{[P1$\rangle$GM]}}
                \texttt{No} \\
\\ 
\texttt{To determine whether having a sister positively affects heart condition *through* blood pressure, we must calculate the **Indirect Effect** (also known as the Mediated Effect) of having a sister on heart condition via blood pressure.} \\
\\ 
\texttt{In causal inference, the indirect effect is the change in the outcome (Heart Condition) that occurs when the treatment (Having a Sister) changes the mediator (Blood Pressure), while holding the treatment status constant at the level of the treated group (or comparing the counterfactuals).} \\
\\ 
\texttt{Let's define the variables:} \\
\texttt{*   \$S\$: Having a sister (\$S=1\$) vs. No sister (\$S=0\$).} \\
\texttt{*   \$B\$: Blood pressure (\$B=1\$ for High, \$B=0\$ for Low).} \\
\texttt{*   \$H\$: Heart condition (\$H=1\$ for Healthy, \$H=0\$ for Unhealthy).} \\
\\ 
\texttt{**1. Extract the Probabilities from the Prompt:**} \\
\texttt{*   \$P(H=1 | S=0, B=0) = 0.63\$} \\
\texttt{*   \$P(H=1 | S=0, B=1) = 0.31\$} \\
\texttt{*   \$P(H=1 | S=1, B=0) = 0.72\$} \\
\texttt{*   \$P(H=1 | S=1, B=1) = 0.41\$} \\
\texttt{*   \$P(B=1 | S=0) = 0.66 \textbackslash{}implies P(B=0 | S=0) = 0.34\$} \\
\texttt{*   \$P(B=1 | S=1) = 0.38 \textbackslash{}implies P(B=0 | S=1) = 0.62\$} \\
\\ 
\texttt{**2. Analyze the Pathways:**} \\
\texttt{The question asks specifically about the effect *through* blood pressure. This isolates the path \$S \textbackslash{}to B \textbackslash{}to H\$.} \\
\texttt{The indirect effect is calculated by comparing the expected outcome if the mediator (\$B\$) were set to the distribution it would have under \$S=1\$ versus \$S=0\$, while keeping the direct effect of \$S\$ on \$H\$ (the path \$S \textbackslash{}to H\$) constant.} \\
\\ 
\texttt{However, a simpler and standard way to interpret "effect through blood pressure" in this context (where we have full conditional probabilities) is to look at the **Natural Indirect Effect (NIE)**.} \\
\texttt{The NIE is defined as:} \\
\texttt{\$NIE = E{[}Y(1, B(1)){]} {-} E{[}Y(1, B(0)){]}\$} \\
\texttt{Where:} \\
\texttt{*   \$Y(1, B(1))\$ is the outcome if everyone has a sister and their blood pressure is whatever it naturally would be if they had a sister.} \\
\texttt{*   \$Y(1, B(0))\$ is the outcome} \\
            }
        }
    }
    & & \\ \\

    \theutterance \stepcounter{utterance}  
    & & & \multicolumn{2}{p{0.3\linewidth}}{
        \cellcolor[rgb]{0.9,0.9,0.9}{
            \makecell[{{p{\linewidth}}}]{
                \texttt{\tiny{[GM$|$GM]}}
                \texttt{No} \\
            }
        }
    }
    & & \\ \\

    \theutterance \stepcounter{utterance}  
    & & & \multicolumn{2}{p{0.3\linewidth}}{
        \cellcolor[rgb]{0.9,0.9,0.9}{
            \makecell[{{p{\linewidth}}}]{
                \texttt{\tiny{[GM$|$GM]}}
                \texttt{yes} \\
            }
        }
    }
    & & \\ \\

    \theutterance \stepcounter{utterance}  
    & & & \multicolumn{2}{p{0.3\linewidth}}{
        \cellcolor[rgb]{0.9,0.9,0.9}{
            \makecell[{{p{\linewidth}}}]{
                \texttt{\tiny{[GM$|$GM]}}
                \texttt{game\_result = LOSE} \\
            }
        }
    }
    & & \\ \\

\end{supertabular}
}

\end{document}
